% answers_short.tex —— 2009 年·普通高等学校招生考试·陕西卷·理科 · 简短答案速查
% 由 tools/make_short_answers.py 自 answers.tex 的 \gksolhead/\gkanswer 提取，
% 并与 2009/2009_数学真题_陕西理_答案要点.md 逐题比对后生成，勿手改（改源 MD 后重跑）。

\begin{gknotice}
\textbf{使用说明}：本节只给最终答案便于快速核对，完整推导见解析版。答案为编写者独立推导，\textbf{非官方标准答案}。
\end{gknotice}

\gkpart{选择题}
\begin{center}\begin{tabular}{ccccccccccccc}
\toprule
题号 & 1 & 2 & 3 & 4 & 5 & 6 & 7 & 8 & 9 & 10 & 11 & 12 \\
\midrule
答案 & A & D & B & D & A & C & C & A & C & B & B & C \\
\bottomrule
\end{tabular}\end{center}

\gkpart{填空题}
\begin{itemize}[leftmargin=2.6em]
  \item[\textbf{13.}] $1$。
  \item[\textbf{14.}] $8$。
  \item[\textbf{15.}] $\dfrac\pi2$。
  \item[\textbf{16.}] $-2$。
\end{itemize}

\gkpart{解答题}
\begin{itemize}[leftmargin=2.6em]
  \item[\textbf{17.}] （1）$f(x)=2\sin\left(2x+\dfrac\pi6\right)$；（2）值域 $[-1,2]$。
  \item[\textbf{18.}] （1）证明见解析；（2）$\arccos\dfrac{\sqrt{15}}5$（正切 $\dfrac{\sqrt6}3$）。
  \item[\textbf{19.}] （1）$a=0.2$，$E\xi=1.7$；（2）$0.17$。
  \item[\textbf{20.}] （1）$a=1$；（2）$a\geqslant2$ 时在 $[0,+\infty)$ 递增；$0<a<2$ 时减区间 $\left[0,\sqrt{\dfrac{2-a}a}\right)$，增区间 $\left(\sqrt{\dfrac{2-a}a},+\infty\right)$；（3）$a\geqslant2$。
  \item[\textbf{21.}] （1）$\dfrac{y^2}4-x^2=1$；（2）$\left[2,\dfrac83\right]$。
  \item[\textbf{22.}] （1）数列不单调；奇数项子列递增、偶数项子列递减，均趋于 $\dfrac{\sqrt5-1}2$；（2）证明见解析。
\end{itemize}
